%% Conclusion — \neven{} Paper

\section{Conclusions and Future Work}
\label{sec:conclusion}

\subsection{Summary}

We have presented \neven{}, an open-source polyglot infrastructure that bridges the gap between spreadsheet accessibility and statistical computing power. The system integrates R~4.4.1, Julia~1.12.6, and Python~3.13 as process-isolated language runtimes communicating with Microsoft Excel through Named Pipes and Protocol Buffers serialization. Key characteristics include:

\begin{itemize}
    \item Over 90 R statistical procedures and 70 Julia mathematical functions accessible as native Excel formulas, requiring no programming expertise from end users.
    \item A security sandbox blocking 30+ dangerous patterns per language with five bypass-prevention mechanisms, validated by 154 automated tests.
    \item Interactive visualization through an embedded WebView2 viewer supporting Plotly, D3.js, Leaflet, and rpivotTable.
    \item An in-process WebView2-based REPL console replacing the legacy Electron application, eliminating 50+ CVEs and all Node.js dependencies while providing interactive R/Julia/Python execution with command history, multi-line input, and inline graphics.
    \item Reactive notebooks (Pluto.jl) with bidirectional data exchange between Excel and Julia.
    \item Reproducible report generation via Quarto integration.
    \item A total of 357 automated tests covering sandbox verification, input sanitization, IPC frame validation, type conversion, RAII semantics, property-based correctness proofs, and end-to-end workflows.
\end{itemize}

The system has been validated against established econometric benchmarks \citep{Wooldridge2016}, demonstrating that the IPC architecture introduces no numerical degradation. An enterprise analytics scenario illustrates the potential for reducing tool fragmentation, eliminating data handoff errors, and providing auditable computational workflows in regulated industries.

\paragraph{Python Resilience Engineering.}
The Python integration (ControlPython.exe) required solving four stability bugs before achieving production readiness: (1)~intermittent \texttt{startup.py} failures resolved through a retry mechanism with \texttt{PyErr\_Clear()} between attempts; (2)~\texttt{STATUS\_STACK\_BUFFER\_OVERRUN} crashes during CPython initialization mitigated by reordering initialization after pipe setup and adding SEH exception guards; (3)~\texttt{IndentationError} from line-by-line startup code transmission resolved by sending the entire script as a single code block for Python (preserving blank lines and docstrings); and (4)~indefinite blocking when Python became unresponsive, resolved by adding health-status checks before IPC calls. This engineering effort demonstrates that process-isolated architectures, while providing fault tolerance, require careful attention to initialization ordering and failure recovery---lessons applicable to any multi-process system.

\subsection{Significance}

\neven{} addresses a real and documented problem: the tension between tool accessibility and computational rigor. By preserving Excel as the user interface while delegating computation to established scientific languages, the system enables domain experts to access advanced statistical methods without crossing the programming expertise barrier. The process-isolated architecture ensures that the reliability expectations of spreadsheet users are maintained---a crash in any language runtime does not affect the host application.

The open-source release under GPL~v3 ensures that the system can be freely adopted in academic settings, where budget constraints often prevent adoption of commercial alternatives such as PyXLL or xlwings PRO. The 357 automated tests and comprehensive documentation (12 user-facing chapters, 95 documented functions with examples) provide a foundation for community contributions and long-term maintenance.

\subsection{Future Work}

We identify four directions for future research and development:

\paragraph{OS-level security sandbox.} The current pattern-based sandbox provides protection against common misuse but cannot guarantee security against determined adversaries. Implementing Windows AppContainer isolation for child processes would provide defense-in-depth, restricting file system and network access at the operating system level regardless of the code being executed. This represents a publishable contribution at the intersection of systems security and scientific computing.

\paragraph{Formal user study.} A controlled study with university students (target: econometrics and statistics courses at the Universidad de Costa Rica) would provide empirical evidence for the usability claims. The study design would compare task completion time, error rates, and user satisfaction between traditional R programming and \neven{} formula-based workflows for equivalent statistical analyses.

\paragraph{Performance benchmarks.} Systematic latency measurements comparing \neven{}'s IPC overhead against in-process alternatives (PyXLL, VBA) and COM-based alternatives (xlwings) would quantify the cost of process isolation. These benchmarks would inform users about the performance envelope and guide optimization efforts for bulk-operation scenarios.

\paragraph{AI-assisted interpretation.} The current Python integration includes preliminary support for LLM-based interpretation of statistical results via \texttt{=P.ai\_call(data, prompt)}, supporting both cloud providers (OpenAI, Azure) and local models (Ollama, LM~Studio) for data privacy. A deeper investigation of how large language models can assist non-expert users in interpreting statistical outputs---while maintaining scientific rigor and avoiding hallucinated conclusions---represents an emerging research direction at the intersection of AI and statistical computing.

\subsection{Availability}

\neven{} is released under the GNU General Public License v3 and is available at:
\begin{itemize}
    \item \textbf{Source code:} \url{https://github.com/minor-bonilla/NEVEN}
    \item \textbf{Documentation:} Included in the repository (\texttt{docs/} directory)
    \item \textbf{Installation:} Automated installer (\texttt{Install-NEVEN.exe}) with runtime detection
    \item \textbf{Examples:} Pre-built Excel workbooks with validated formulas (\texttt{Ejemplos/} directory)
\end{itemize}

The system requires Windows~10 or later (64-bit), Microsoft Excel~2013 or later, and at least one of: R~4.4.1+, Julia~1.12.6+, or Python~3.13+. The Julia sysimage (\textasciitilde415~MB) is optional but recommended for eliminating cold-start latency.
